\section{Controlled Mechanism Study}
\label{sec:study}

The case study located the authority transition in an independently published
lifecycle; this section isolates its mechanism in
Safe-Control-Gym~\cite{yuan_safe_control_gym_2022}.  A prospective
contract-separation audit first tests whether finite release evidence supports
both permanent raw release and resident predictive authority, with and without
the attacker, using a one-step resident kernel as a negative control and a
horizon sweep to bound the required lookahead.  Cross-system, detectability, and
runtime-cost studies then characterize the scope and operational consequences
of the fixed attack.  Appendix~\ref{sec:appendix-benchmarks} reports the
low-dimensional semantic benchmarks that separate current-action filtering
from update admission.

\subsection{Prospective Contract-Separation Audit}
\label{sec:release-contract-audit}

All results in this subsection assume the reward-ingestion interface in
Table~\ref{tab:trust-boundary}.  Development grids selected target $(18,-5)$
using five-step acceptance and 120-step harm.  Before the sequential V3 audit,
we fixed the target, learner, state pools, stop criterion, and bounded
perturbation $2\tanh(\epsilon_t(\mu^\star_t-\mu_t))$ on the development run
indexed by seed 2070; runs 2071 and 2072 were then evaluated under the
unchanged protocol.

All 72 poisoned snapshot--state pairs pass the initial five-step check.
Poisoned permanent raw release then failed on 11, 5, and 7 pairs across the
three training runs, versus paired clean counts 0, 0, and 2: 21 of 23 poisoned
failures are poison-only, two shared, none clean-only (two-sided exact paired
$p=9.537\times10^{-7}$), and the CasADi audit identifies the same 23 failures.
Resident predictive authority incurs 0/72 failures and switches four to five
steps before each corresponding poisoned raw-release failure.

The frozen protocol was then evaluated once on five new learner/evaluation
pairs.  The target, attacker, learner, budget, state selection, and contracts
remained fixed.  All 120 snapshot--state pairs passed the initial five-step
check.  Poisoned raw release failed on 5, 4, 7, 2, and 9 pairs across the five
training runs; the paired clean counts were 2, 0, 0, 0, and 0.  Resident
predictive authority incurred 0/120 failures and switched before all 27
poisoned failures.  The 25 poison-only and zero clean-only discordances gave a
two-sided exact $p=5.960\times10^{-8}$.  Poisoned failure occurred in each of
the five independently trained runs.

\emph{The gap without an attacker.}  Two of the 120 clean pairs failed under
permanent raw release after passing the same five-step check that accepted
them.  Two in 120 is a small rate, and the paired design was built to detect
the poisoning effect rather than to estimate the clean one, so we do not claim
that unattacked release is unsafe in general.  We claim what the two failures
witness: acceptance and failure are separated here only by horizon, with no
adversary in between, and the reward-record channel widens a gap that is
already present.

State--snapshot pairs that share a trained learner are not independent
training replicates, so the V3 and V4 exact tests summarize within-snapshot
discordance; replication across training realizations is reported through the
three V3 and five V4 runs.  For independent replication and the one-step
negative control, snapshots from run indices 2040--2042 used disjoint
evaluation pools 8040--8042; 71 of 72 states passed the initial check, and all
mechanisms shared state--snapshot keys.

\begin{table}[t]
\centering
\caption{Contract separation after initial five-step acceptance under bounded
reward-record write access.  V4 is the five-run confirmation; V3 contains the
development and sequential audits.  The independent audit includes the
one-step negative control.  A switch transfers control to the trusted LQR.}
\label{tab:release-contract}
\small
\begin{tabular}{@{}p{0.46\columnwidth}p{0.13\columnwidth}p{0.25\columnwidth}@{}}
\toprule
Contract & Failures & Authority event\\
\midrule
\multicolumn{3}{@{}l}{\emph{V4 confirmation (120 pairs; 5 training runs)}}\\
Clean permanent raw release & 2/120 & ---\\
Poisoned permanent raw release & 27/120 & ---\\
Resident predictive authority & 0/120 & 27 switches\\
\addlinespace
\multicolumn{3}{@{}l}{\emph{Development/sequential V3 (72 paired states)}}\\
Clean permanent raw release & 2/72 & ---\\
Poisoned permanent raw release & 23/72 & ---\\
Resident predictive authority & 0/72 & 23 switches\\
\addlinespace
\multicolumn{3}{@{}l}{\emph{Independent audit (71 accepted states)}}\\
Permanent raw release & 19/71 & ---\\
One-step resident kernel & 14/71 & 79 int.; 41 empty\\
Resident predictive authority & 0/71 & 19 switches\\
\bottomrule
\end{tabular}
\end{table}

Table~\ref{tab:release-contract} makes the three evidence blocks and
deployment contracts explicit.

\emph{How much lookahead a retained check needs.}  A locked sweep varies only
the retained predictor's horizon on the pooled V3 cohort of 72 pairs, holding
attack, snapshots, and deployment states fixed.  Three quantities move
together.  Release admission falls as the horizon grows: 72, 72, 72, 57, and
49 of 72 candidates accepted at horizons 1, 3, 5, 10, and 20.  Switching lead
time grows: the median switch preceded the corresponding raw-release failure
by 0, 2, 4, and 9 steps at horizons 1, 3, 5, and 10.  Resident failures track
lead time rather than residence: 12 at horizon 1, where switches occurred with
zero lead, and 0 at horizons 3, 5, and 10.  Horizon 1 failed the predeclared
separation condition for insufficient lookahead; horizon 20 failed it from the
opposite side, rejecting so many candidates that none of the 49 admitted pairs
retained a poisoned failure to separate.  The usable lookahead is an interval,
and its lower edge is set by the lead time the trusted baseline needs, not by
the accuracy of the prediction.  Figure~\ref{fig:horizon-contract} plots the
three curves.  Shielding work already observes that a finite-horizon shield
needs enough lookahead or the agent reaches states where every action is
unsafe~\cite{konighofer_online_shielding_2022}; the upper edge measured here
is a different cost, paid in release admission rather than in computation.

\begin{figure}[t]
\centering
\includegraphics[width=0.99\columnwidth]{fig10_horizon_contract.pdf}
\caption{The retained check's horizon defines a usable interval rather than a
threshold.  Pooled V3 cohort, 72 snapshot--state pairs, with attack,
snapshots, and deployment states fixed; shading marks the horizons meeting the
predeclared separation condition.  At horizon 1 the retained check switches
with zero lead and still leaves 12 failures; at horizon 20 the release check
rejects so many candidates that no poisoned failure remains to separate.}
\label{fig:horizon-contract}
\end{figure}

In the independent audit, permanent raw release failed on 6/24,
3/23, and 10/24 accepted pairs.  Resident predictive authority repeats the
same five-step raw-policy prediction at every physical step and switches on
exactly those 19 paired states, four steps before failure; no physical failure
remained.  In contrast, the one-step resident kernel failed on 14 pairs despite
79 interventions and reached an empty kernel 41 times.  Repeated five-step
prediction therefore provided an effective recoverability trigger; one-time
authorization and one-step filtering did not.

\subsection{Reward-Influence Geometry and Optimizer Boundary}

We validated the reward-influence analysis on an unclipped cartpole batch
from the learner indexed by seed 2040.  Reconstructing the implemented
gradient from Equation~\eqref{eq:normalized-update} has error
$2.78\times10^{-17}$, and across 22 release-gate halfspaces the largest
discrepancy between an SOCP witness and its computed support is
$1.07\times10^{-9}$.  Audited on all 36 poisoned V3 batches, the homogenized
clipped formulation yields verified support witnesses on 22/36, where the
earlier premises admitted 0/36.  Exact batch support nevertheless did not
explain the successful multi-batch attack: the physical experiment evaluates
the fixed target and bounded attacker without establishing joint gate
optimality of the \(\tanh\) rule.  Appendix~\ref{sec:appendix-influence}
reports the full implementation audit and its parameter-clipping boundary.

\subsection{Cross-System Reward-Record Poisoning}

We next tested the lifecycle failure with updates produced by an online
learner: a Gaussian residual actor on a trusted controller, with standard
reward-to-go REINFORCE gradients computed from PyBullet transitions.  A backup
action filter protects the adaptation trajectory, and deployment then audits
the raw learned snapshot, the object requiring lifecycle admission.  Training,
certificate, calibration, and deployment states used disjoint run indices; all
mechanisms within a training realization shared paired state sets; and attack
directions and hyperparameters were fixed on a development realization before
the three-run audits.  The exact tests operate at the paired
deployment-state-block level conditional on the three learned snapshots, so
their $p$-values summarize within-snapshot discordance, and Wilson and
paired-block bootstrap intervals have the same conditional scope.

\begin{table*}[t]
\centering
\caption{Paired reward-record poisoning and snapshot admission. ``Commit''
denotes finite-horizon admission; permanent mechanisms use the common-state
audit in Fig.~\ref{fig:unified-frontier}. Raw-snapshot rewards after failure
are descriptive, not utility comparisons.}
\label{tab:end-to-end}
% Auto-generated by generate_tdsc_result_table.py; do not edit.
\begin{tabular}{llrrrr}
\toprule
System & Mechanism & Learner seeds & Rollouts & Violations & Mean reward \\
\midrule
Cartpole & Clean REINFORCE & 3 & 188 & 1 & $-0.197$ \\
Cartpole & Poisoned action-only & 3 & 188 & 38 & $-1.269$ \\
Cartpole & Poisoned commit & 3 & 188 & 0 & $-0.102$ \\
Cartpole & Always-freeze & 3 & 188 & 0 & $-0.117$ \\
\midrule
Quadrotor & Clean REINFORCE & 3 & 72 & 0 & $-0.030$ \\
Quadrotor & Poisoned action-only & 3 & 72 & 68 & $-0.170$ \\
Quadrotor & Poisoned commit & 3 & 72 & 0 & $-0.060$ \\
Quadrotor & Always-freeze & 3 & 72 & 0 & $-0.029$ \\
\bottomrule
\end{tabular}

\end{table*}

Table~\ref{tab:end-to-end} reports the paired outcomes for both systems.
On cartpole, the clean learner violates 1/188 admitted rollouts and the
poisoned action-only learner violates 38/188.  There are 37 paired poison-only
failures and no clean-only failures (two-sided exact paired
$p=1.455\times 10^{-11}$).  Both commit admission and always-freeze have zero
observed violations.

On two-dimensional quadrotor, clean learning has 0/72 violations while
reward-only poisoning causes 68/72, all delayed, across all three training
realizations (68 poison-only, zero clean-only, two-sided exact
$p=6.776\times10^{-21}$).  Every one of those failures is delayed past
adaptation, so protected execution reports nothing about any of them: the
separation follows the authority transition rather than the platform.  Commit
admission and always-freeze again have zero observed violations, and the gate
retains $0.55$, $0.40$, and $0.70$ of the poisoned pending updates, which
establishes update availability under attack rather than benign learning
utility.

\emph{Learned-baseline variant.}  Replacing REINFORCE's constant batch-mean
baseline with a learned per-state linear value baseline, under a protocol
fixed before execution and retaining the same Monte Carlo return, residual
actor, score function, and reward-record channel, poisoning caused 40/126
raw-release failures across three independent training runs, against 2/126 for
paired clean release, with commit admission and always-freeze at 0/126 (38
poison-only, zero clean-only, two-sided exact $p=7.3\times10^{-12}$).  This is
a robustness check within the Monte Carlo policy-gradient family, not a
distinct learner family: off-policy and bootstrapped learners change the
attack channel through a temporal-difference target, so the official-PPO audit
below remains the learner boundary.  These rates use the full 42-state
envelope ($126=3\times42$) without initial-acceptance filtering and are not
directly comparable to the filtered 120-pair confirmation cohort.

\emph{Non-singleton adaptation set.}  Adding the observation-FDI channel of
Section~\ref{sec:threat} to the reward-record attacker (a pole-angle bias of
radius $\rho$, so the protective kernel ranges over a box of half-width
$\rho$, about $\pm2.3^\circ$ at $\rho=0.04$) left the release certificate and
deployment states unchanged and preserved the contract separation at every
non-singleton radius.  Poisoned raw release failed on 44/126 and 61/126 pairs
at $\rho=0.02$ and $0.04$, against 4/126 and 3/126 for clean release, with
commit admission and always-freeze at 0/126; the near-singleton reference at
$\rho=0.005$ was 41/126 versus 5/126.  Failure counts were non-monotone across
training realizations, while the kernel remained nonempty throughout the grid.

\subsection{Reward-Record Integrity and Detectability Boundary}

We evaluated attack detectability separately from deployment failure.  Four
checks were fixed on the clean development run indexed by seed 2070.  The
frozen checks were then evaluated on 84 clean and 84 poisoned batches from
seven independent training runs.

\begin{table}[t]
\centering
\caption{Frozen reward-integrity audit.  Rates are batch-level true- and
false-positive rates (TPR/FPR) across seven training runs.  Each check requires
the stated defender assumption.}
\label{tab:reward-integrity}
\footnotesize
\begin{tabular}{@{}p{0.27\columnwidth}p{0.43\columnwidth}p{0.17\columnwidth}@{}}
\toprule
Check & Required trusted information & TPR/FPR\\
\midrule
Reward recomputation & Authenticated transition and deterministic reward
& 1.000/0\\
Known-sign check & Task guarantee that reward is nonpositive & 0.905/0\\
Scalar envelope & Clean development range transfers & 0.917/0.036\\
Batch-mean envelope & Clean batch mean transfers & 0.012/0.036\\
\bottomrule
\end{tabular}
\end{table}

Table~\ref{tab:reward-integrity} shows that trusted recomputation blocks every
evaluated edit, and even task-specific scalar checks detect over 90\% of
poisoned batches.  The ineffective batch-mean envelope shows why a generic
``monitor the average'' statement is not enough: detection depends on reward
semantics and trusted inputs.  Cryptographic origin binding would also reject
post-producer mutation if the trusted producer signs the reward together with
its transition, action, and time.  That producer remains within the trusted
computing base of the evaluated threat model.

Temporal and moment-preserving edits tested on the development realization
satisfied their integrity constraints and moved parameters toward the harmful
target, yet each caused 0/24 raw-release failures.  The effective fixed attack
was therefore readily detectable by the task-aware checks.  Authentication and
trusted reward recomputation provide upstream controls; commit admission and
resident predictive authority provide downstream containment.

\emph{In-loop known-sign gate.}  We placed the same known-sign check inside the
update loop and froze batches that violated the nonpositive reward guarantee.
Across three independent training runs, this gate reduced poisoned raw-release
failures from 39/126 to 0/126 and froze 35 of 36 poisoned batches.  The paired
defended-versus-undefended test gave $p=3.6\times10^{-12}$.  The gate froze no
clean batch but approached always-freeze availability under attack.  It detects
the sign-violating effective attack, while the sign-respecting variants evade
the check and independently caused no deployment harm.  This gate implements
the measured detectability boundary.

\subsection{Runtime Escape Conditions and Their Cost}

Permanent runtime filtering provides a second architectural response.  Its
certificate covers the composed policy $F\circ\pi$, allowing the filter to
contain raw-policy failure before physical execution.  We compared permanent
and snapshot contracts on a result-independent pool of 144 states from three
disjoint run indices.  For each training realization the common operating set
held states certified for both freeze and the committed snapshot at horizon
100 and guard 0.003, and eight order-spanning states were evaluated from
identical state/reset pairs, giving each mechanism 24 paired rollouts and 2400
possible steps.


\begin{figure*}[t]
\centering
\includegraphics[width=0.94\textwidth]{tdsc_security_availability_frontier.pdf}
\caption{The unified frontier on the same paired state blocks.  Error bars in
(a) are 95\% Wilson intervals for violation rate and paired-block bootstrap
intervals for mean rollout reward, conditional on the three locked learner
snapshots.  Panel (b) separates the negligible raw
policy latency from the permanently-online protection cost; raw, commit, and
freeze mechanisms overlap near 0.026 ms and zero intervention.}
\label{fig:unified-frontier}
\end{figure*}

Table~\ref{tab:unified-frontier} in
Appendix~\ref{sec:appendix-frontier} and Fig.~\ref{fig:unified-frontier} show
the resulting separation.  The poisoned
raw snapshots violate 23/24 rollouts and complete only one.  Clean, commit,
freeze, the permanent five-step backup shield, and Safe-Control-Gym's
official linear model predictive safety certification (MPSC) mechanism all
have zero observed failures and complete 24/24.
The shield intervenes on 511/2400 steps with no rejection, while MPSC
intervenes on 2400/2400 with no infeasible solve, at median online latencies
of 222 and 60 ms against 0.026 ms for raw snapshot evaluation.  These are
single-process wall-clock measurements, not worst-case execution-time
guarantees.

The safe mechanisms differ in utility.  Mean reward ordered freeze, clean,
commit, shield, MPSC from $-0.023$ to $-0.457$; every paired rollout delta
relative to freeze had a bootstrap 95\% interval excluding zero, and no
mechanism improved on freeze in any of the 24 blocks.  These values
characterize the attack-evaluation frontier; the benign adaptation audit below
evaluates learning utility separately.



Table~\ref{tab:commit-availability} in
Appendix~\ref{sec:appendix-boundaries} makes the snapshot-admission cost
explicit: common operating-set coverage is 99.3\%, 100.0\%, and 90.3\% across
the three training runs, and trusted-state admission took 8.49--9.32\,s per
run.  Offline commit therefore avoids permanent per-action computation but
pays finite certificate coverage and update latency.  The comparison distinguishes three
contracts: certify a snapshot before raw deployment, retain a filter
permanently, or stop adapting.

\subsection{Sizing a Release Gate}

The measurements above support a contract-specific rule.  Two decisions are
made at the authority transition, and each has a measurable input.

\emph{Retiring authority.}  The release evidence must re-establish every
premise the transition changes (Proposition~\ref{prop:transfer}).  Here those
premises are authority and horizon: the accepted check bounds the raw policy
for five steps, and only while it can still hand control back.  Re-running the
same check discharges neither obligation, which is why permanent raw release
failed on states it had accepted, with and without an attacker.

\emph{Retaining authority.}  The lookahead then has a lower and an upper edge,
and both are measurable on the deploying system without knowing the attack.
The lower edge is the lead time the trusted baseline needs to recover: at
horizon 1 the retained check switched with zero lead and still left 12
failures, whereas horizons 3 to 10 switched 2 to 9 steps early and left none.
The upper edge is the admission the operator can afford to lose: raising the
horizon from 5 to 10 and to 20 rejected 15 and 23 of the 72 candidates the
shorter checks admitted.  Both edges come from the sweep of
Fig.~\ref{fig:horizon-contract} applied to the deployer's own baseline and
candidate pool, so the gate can be sized from the recovery dynamics of the
backup rather than from an attack model.

\subsection{Boundaries and Validity}

Four locked audits bound the result.  A trusted state anchor shrinks
certificate uncertainty without sanitizing the reward record: every anchor
period retains all 16 certificate centers, has 0/24 violations on new paired
blocks, and improves no reward over freeze.  Against the official pretrained
PPO policy, which is safe on 72/72 held-out rollouts, the locked
$L_\infty\le0.5$ budget produced 1/24 poisoned failures at one checkpoint and
0/24 at another, so the audit showed checkpoint-level update redirection
without a stable deep-PPO unsafe-policy effect.  Under a benign $0.003$\,N
differential actuator bias, freeze, clean adaptation, and the
certificate-gated learner were all violation-free; the gate accepted a nonzero
update in 11/12 batches but improves no paired rollout over freeze, so
certificate acceptance did not establish useful adaptation.  A fourth audit bounds how far the finite certificate transfers.  On 1728
snapshot--state pairs at the pre-existing 100-step, 0.003-guard configuration,
CasADi certified 1293, none of which violated the PyBullet state box within
100 steps, while the poisoned raw snapshots violated 403/432 pairs at
certificate coverage 0.062 --- informative rejection rather than vacuous
evidence.  The guard margin is what separates the two regimes:
without it the poisoned snapshots produced 1, 2, and 1 false acceptances at
horizons 20, 50, and 100, and every locked positive guard removed the observed
false acceptances while rejecting more physically safe pairs
(Appendix~\ref{sec:appendix-coverage}).  The gate therefore has a second
measurable knob beyond horizon, calibrated against model mismatch rather than
against an attacker.

The CasADi gates use fixed finite state
sets while physical outcomes use independent PyBullet rollouts; they are
sampled finite-horizon certificates, not continuous-state invariant-set
proofs.  Realized state violations and pre-clip saturation are reported
separately.  Appendix~\ref{sec:appendix-boundaries} gives the full figures.
